\documentclass[letterpaper]{article}
\usepackage{common/ohpc-doc}
\setcounter{secnumdepth}{5}
\setcounter{tocdepth}{5}

% Include git variables
\input{vc.tex}

% Define Base OS and other local macros
\newcommand{\OS}{Rocky}
\newcommand{\baseOS}{\OS{[} 9.6}
\newcommand{\OSRepo}{\OS{}\_9.6}
\newcommand{\OSTree}{EL\_9}
\newcommand{\OSTag}{el9}
\newcommand{\baseos}{rocky9.6}
\newcommand{\baseosshort}{rocky9}
\newcommand{\provisioner}{OpenCHAMI}
\newcommand{\provheader}{\provisioner{}}
\newcommand{\rms}{SLURM}
\newcommand{\rmsshort}{slurm}
\newcommand{\arch}{x86\_64}

\newcommand{\isoimage}{Rocky-9.4-x86\_64-dvd.iso}
\newcommand{\downloadpage}{https://rockylinux.org/download/}
\newcommand{\confluentprofile}{rocky-9.4-x86\_64-default}
\newcommand{\OpenCHAMI}{OpenCHAMI}
\newcommand{\VERLONG}{0.0.31}
% Define package manager commands
\newcommand{\pkgmgr}{dnf}
\newcommand{\addrepo}{dnf -y config-manager --add-repo}
\newcommand{\chrootaddrepo}{nodeshell compute dnf -y config-manager --add-repo}
\newcommand{\clean}{dnf clean expire-cache}
\newcommand{\chrootclean}{dnf --installroot=\$CHROOT clean expire-cache}
\newcommand{\install}{dnf -y install}
\newcommand{\chrootinstall}{dnf -y --installroot=\$CHROOT install}
\newcommand{\groupinstall}{dnf -y groupinstall}
\newcommand{\groupchrootinstall}{dnf -y --installroot=\$CHROOT groupinstall}
\newcommand{\rpminstall}{rpm}
\newcommand{\remove}{dnf -y remove}
\newcommand{\upgrade}{dnf -y upgrade}
\newcommand{\chrootupgrade}{dnf -y --installroot=\$CHROOT upgrade}
\newcommand{\tftppkg}{syslinux-tftpboot}
\newcommand{\beegfsrepo}{https://www.beegfs.io/release/beegfs\_7.2.1/dists/beegfs-rhel8.repo}
\newcommand{\cudarepo}{https://developer.download.nvidia.com/compute/cuda/repos/rhel9/x86\_64/cuda-rhel9.repo}
% boolean for os-specific formatting
\toggletrue{isCentOS}
\toggletrue{isCentOS_ww_slurm_x86}
\toggletrue{isSLURM}
\toggletrue{isx86}
\toggletrue{isCentOS_x86}

\begin{document}
\graphicspath{{common/figures/}}
\thispagestyle{empty}

% Title Page
\input{common/title}
% Disclaimer
\input{common/legal}

\newpage
\tableofcontents
\newpage

% Introduction  --------------------------------------------------

\section{Introduction} \label{sec:introduction}
\input{common/install_header}
\input{common/intro} \\

\input{common/base_edition/edition}
\input{common/audience}
\input{common/requirements}
\input{common/inputs}



% Base Operating System --------------------------------------------

\section{Install Base Operating System (BOS)}
\input{common/bos}

%\clearpage
% begin_ohpc_run
% ohpc_validation_newline
% ohpc_validation_comment Disable firewall
\begin{lstlisting}[language=bash,keywords={}]
[sms](*\#*) systemctl disable firewalld || true
[sms](*\#*) systemctl stop firewalld || true
\end{lstlisting}
% end_ohpc_run

% ------------------------------------------------------------------

\section{Install \OpenCHAMI{} and Provision Nodes with SMD}\label{sec:provision_compute_bos}
Installation is completed in a few parts. First we build the images that the
{\em compute} nodes will boot off of. Then we configure the cloud-init steps
to add a local account, mount remote file systems, and inject the munge key
into the OS. \OHPC{} components will be added to both the SMS and compute nodes.
The SMS will get the components during install, and compute nodes will get them
during image builds.


\subsection{\OpenCHAMI{} prerequisites}\label{sec:openchami_prerequisites}
Add SMS name to /etc/hosts so we can make API calls by hostname
% begin_ohpc_run
% ohpc_validation_newline
% ohpc_validation_comment
\begin{lstlisting}[language=bash,keywords={}]
[sms](*\#*) echo "${sms_ip} ${sms_name}" >> /etc/hosts
\end{lstlisting}
% end_ohpc_run


\subsection{Enable \OpenCHAMI{} repository for local use}\label{sec:enable_openchami}
To begin we need to add the OpenCHAMI RPM to the SMS node, by adding it to the list of
available packages locally. This requires network access from the {\em master}
node to the internet.

% begin_ohpc_run
% ohpc_validation_newline
% ohpc_comment_header Enable OpenCHAMI Repositories and Support \ref{sec:enable_openchami}
\begin{lstlisting}[language=bash,keywords={},basicstyle=\fontencoding{T1}\fontsize{8.0}{10}\ttfamily,
	literate={VER}{\OHPCVerTree{}}1 {OSREPO}{\OSTree{}}1 {TAG}{\OSTag{}}1 {ARCH}{\arch{}}1 {-}{-}1]i
[sms](*\#*) (*\install*) podman buildah ansible-core nfs-utils tcpdump tftp curl yq
[sms](*\#*) rpm --import https://github.com/OpenCHAMI/release/releases/download/vVERLONG/public_gpg_key.asc
[sms](*\#*) (*\install*) https://github.com/OpenCHAMI/release/releases/download/vVERLONG/openchami-VERLONG.rpm
\end{lstlisting}
% ohpc_validation_newline
% end_ohpc_run


\subsection{Setting no proxy}\label{sec:openchami_noproxy}

\OpenCHAMI{} runs its own internal proxy for services which may not interact well
if you run another httpd server. Here we set the no proxy variable to prevent lookups.
% begin_ohpc_run
% ohpc_comment_header OpenCHAMI No Proxy Settings \ref{sec:openchami_noproxy}
\begin{lstlisting}[language=bash,keywords={}]
[sms](*\#*) export additonal_no_proxy="${sms_name},opaal,opaal-idp,smd,smd-init,hydra,cloud-init,acme-deploy,bss,haproxy,postgres,step-ca,configurator"
[sms](*\#*) if [ -z "$no_proxy" ]; then
[sms](*\#*)     export no_proxy="$additonal_no_proxy"
[sms](*\#*) else
[sms](*\#*)     export no_proxy="$no_proxy,$additonal_no_proxy"
[sms](*\#*) fi
[sms](*\#*) if [ -z "$NO_PROXY" ]; then
[sms](*\#*)     export NO_PROXY="$additonal_no_proxy"
[sms](*\#*) else
[sms](*\#*)     export NO_PROXY="$NO_PROXY,$additonal_no_proxy"
[sms](*\#*) fi

\end{lstlisting}
% ohpc_validation_newline
% end_ohpc_run


\subsection{Setup DHCP}\label{sec:dhcp_config}
% -*- mode: latex; fill-column: 120; -*-
OpenCHAMI uses CoreDHCP as the DHCP service to provide IP Addresses to nodes,
there is a default configuration in place that we need to overwrite.

Of Note if your router address is different than the internal network address .254 you will need to change this file on your own.
CoreDHCP has been configured to only provide addresses to nodes it knows.
There are some extra \href{https://github.com/OpenCHAMI/deployment-recipes/blob/main/quickstart/DHCP.md}{configuration changes} you can make to change that

% begin_ohpc_run
% ohpc_comment_header Configure DHCP \ref{sec:dhcp_config}
\begin{lstlisting}[language=bash,keywords={}]
[sms](*\#*) cat > /etc/openchami/configs/coredhcp.yaml <<EOF
[sms](*\#*) server4:
[sms](*\#*)  plugins:
[sms](*\#*)     - server_id: ${sms_ip}
[sms](*\#*)     - router: ${ipv4_gateway}
[sms](*\#*)     - dns: ${dns_servers}
[sms](*\#*)     - netmask: ${internal_netmask}
[sms](*\#*)     - coresmd: https://${sms_name}:8443 http://${sms_ip}:8081 /root_ca/root_ca.crt 30s 1h false
[sms](*\#*) EOF
\end{lstlisting}
% ohpc_validation_newline
% end_ohpc_run


\vspace*{-0.15cm}
\subsection{Complete basic \OpenCHAMI{} setup for {\em master} node}\label{sec:setup_openchami}
At this point, all of the packages necessary to use
\OpenCHAMI{} on the {\em master} host should be installed. Next, we create
and enable the Container Registry and S3 Object store.
These are optional if you have an existing S3 or Container Registry available for use at your local site.

\begin{center}
\begin{tcolorbox}[]
\small The S3 and registry containers are configured as \href{https://docs.podman.io/en/latest/markdown/podman-systemd.unit.5.html}{Podman Quadlets},
which interface with systemd allowing you to perform normal systemd operations on them like:
\begin{lstlisting}[language=bash]
systemctl status registry
systemctl restart minio-server
journalctl -u registry
\end{lstlisting}
\end{tcolorbox}
\end{center}

First create local directories for container storage
% begin_ohpc_run
% ohpc_comment_header Create directories for registry and S3 containers
%\begin{verbatim}
\begin{lstlisting}[language=bash,literate={-}{-}1,keywords={},upquote=true,keepspaces]
[sms](*\#*) mkdir -p /opt/ohpc/admin/data/oci
[sms](*\#*) mkdir -p /opt/ohpc/admin/data/s3
\end{lstlisting}
%\end{verbatim}
% end_ohpc_run

Now create the container registry quadlet definition
% begin_ohpc_run
% ohpc_validation_newline
% ohpc_comment_header Create container registry quadlet
\begin{lstlisting}[language=bash,literate={-}{-}1,keywords={},upquote=true,keepspaces]
[sms](*\#*) cat > /etc/containers/systemd/registry.container <<EOF
[sms](*\#*) # registry.container
[sms](*\#*) [Unit]
[sms](*\#*) Description=Image OCI Registry
[sms](*\#*) After=network-online.target
[sms](*\#*) Requires=network-online.target
[sms](*\#*)
[sms](*\#*) [Container]
[sms](*\#*) ContainerName=registry
[sms](*\#*) HostName=registry
[sms](*\#*) Image=docker.io/library/registry:latest
[sms](*\#*) Volume=/opt/ohpc/admin/data/oci:/var/lib/registry:Z
[sms](*\#*) PublishPort=5000:5000

[sms](*\#*) [Service]
[sms](*\#*) TimeoutStartSec=0
[sms](*\#*) Restart=always
[sms](*\#*)
[sms](*\#*) [Install]
[sms](*\#*) WantedBy=multi-user.target
[sms](*\#*) EOF
\end{lstlisting}
% ohpc_validation_newline
% end_ohpc_run

Then create the local S3 object store quadlet definition
% begin_ohpc_run
% ohpc_validation_newline
% ohpc_comment_header Create Minio S3 object storage quadlet
\begin{lstlisting}[language=bash,literate={-}{-}1,keywords={},upquote=true,keepspaces]
[sms](*\#*) cat > /etc/containers/systemd/minio.container <<EOF
[sms](*\#*) [Unit]
[sms](*\#*) Description=Minio S3
[sms](*\#*) After=local-fs.target network-online.target
[sms](*\#*) Wants=local-fs.target network-online.target
[sms](*\#*)
[sms](*\#*) [Container]
[sms](*\#*) ContainerName=minio-server
[sms](*\#*) Image=docker.io/minio/minio:latest
[sms](*\#*) # Volumes
[sms](*\#*) Volume=/opt/ohpc/admin/data/s3:/data:Z

[sms](*\#*) # Ports
[sms](*\#*) PublishPort=9000:9000
[sms](*\#*) PublishPort=9091:9001

[sms](*\#*) # Environment Variables
[sms](*\#*) Environment=MINIO_ROOT_USER=admin
[sms](*\#*) Environment=MINIO_ROOT_PASSWORD=password1234
[sms](*\#*)
[sms](*\#*) # Command to run in container
[sms](*\#*) Exec=server /data --console-address :9001
[sms](*\#*)
[sms](*\#*) [Service]
[sms](*\#*) Restart=always

[sms](*\#*) [Install]
[sms](*\#*) WantedBy=multi-user.target
[sms](*\#*) EOF
\end{lstlisting}
% ohpc_validation_newline
% end_ohpc_run

We need to reload the system daemon, and to allow the MinIO,
and OCI Registry containers to start.

% begin_ohpc_run
% ohpc_validation_newline
% ohpc_validation_comment Reload Systemd daemon and start the s3 (minio) and container registry
\begin{lstlisting}[language=bash,keywords={},upquote=true,basicstyle=\footnotesize\ttfamily,literate={BOSVER}{\baseos{}}1]
[sms](*\#*) systemctl daemon-reload
[sms](*\#*) systemctl start minio.service
[sms](*\#*) systemctl start registry.service
\end{lstlisting}
% ohpc_validation_newline
% end_ohpc_run


\subsection{Certificate Creation}\label{sec:certificate_creation}
% -*- mode: latex; fill-column: 120; -*-
With \OpenCHAMI{} a minimal open source certificate authority from smallstep is needed. The included automation initialized the CA on first startup. We can immediately download a certificate into the system trust bundle on the host for trusting all subsequent OpenCHAMI certificates. Notably, the certificate authority features ACME for automatic certificate rotation.

% begin_ohpc_run
% ohpc_comment_header Creating OpenCHAMI certificates \ref{sec:certificate_creation}
\begin{lstlisting}[language=bash,keywords={}]
[sms](*\#*) openchami-certificate-update update ${sms_name}
\end{lstlisting}
% ohpc_validation_newline
% end_ohpc_run


\subsection{Starting \OpenCHAMI{}}\label{sec:starting_openchami}
% -*- mode: latex; fill-column: 120; -*-
Starting OpenCHAMI can take a while especially when the node images are being downloaded for the first time.
The of all OpenCHAMI services can be checked with
\begin{lstlisting}[language=bash,keywords={}]
[sms](*\#*) systemctl list-dependencies openchami.target
\end{lstlisting}

% begin_ohpc_run
% ohpc_comment_header Starting OpenCHAMI \ref{sec:starting_openchami}
\begin{lstlisting}[language=bash,keywords={}]
[sms](*\#*) systemctl start openchami.target
\end{lstlisting}
% ohpc_validation_newline
% end_ohpc_run


\subsection{\OpenCHAMI{} Client Install}\label{sec:openchami_client_install}
% -*- mode: latex; fill-column: 120; -*-
\OpenCHAMI{} is very configurable, and it has a client that allows you do most of the configuration
via the command line, so installing that client makes management much easier.
% begin_ohpc_run
% ohpc_comment_header OpenCHAMI Client Install \ref{sec:openchami_client_install}
\begin{lstlisting}[language=bash,keywords={}]
[sms](*\#*) dnf install -y \
	https://github.com/OpenCHAMI/ochami/releases/download/v0.3.4/ochami_0.3.4_linux_amd64.rpm
[sms](*\#*) yes | ochami config cluster set --system --default \
	${compute_prefix} cluster.uri https://${sms_name}:8443
\end{lstlisting}
% ohpc_validation_newline
% end_ohpc_run


% begin_ohpc_run
% ohpc_validation_newline
% ohpc_comment_header Restart opaal and opaal-idp (FIXME: this should not be necessary)
% ohpc_validation_newline
\begin{lstlisting}[language=bash,keywords={}]
[sms](*\#*) systemctl restart opaal-idp opaal
[sms](*\#*) sed -e "s,FQDN|md5sum,RANDOM|sha512sum,g" -i /usr/libexec/openchami/ohpc-nodes.sh
\end{lstlisting}
% ohpc_validation_newline
% end_ohpc_run

\subsection{Node Discovery}\label{sec:static_node_discovery}
% -*- mode: latex; fill-column: 120; -*-
\OpenCHAMI{} can do dynamic discovery, but here we will do node discovery by YAML file.
The YAML file can be imported by the discover command creating a way for adding nodes that
do not have redfish capabilities, or when you just want to provide a list of nodes.


Create the nodes.yaml file the static discovery will use to populate SMD
% begin_ohpc_run
% ohpc_comment_header Discovering Nodes \ref{sec:static_node_discovery}
\begin{lstlisting}[language=bash,keywords={}]
[sms](*\#*) mkdir -p /opt/ohpc/admin/nodes
[sms](*\#*) export ${compute_prefix^^}_ACCESS_TOKEN=$(bash -lc 'gen_access_token')
[sms](*\#*) echo "nodes:" > /opt/ohpc/admin/nodes/nodes.yaml
[sms](*\#*) for((i=0; i < $num_computes; i++)); do
		/usr/libexec/openchami/ohpc-nodes.sh ${c_name[$i]} $i ${c_bmc[$i]} ${c_mac[$i]} ${c_ip[$i]}
       done
\end{lstlisting}
% ohpc_validation_newline
% end_ohpc_run


Now use the generated file to populate the SMD database
% begin_ohpc_run
% ohpc_comment_header Passing ochami the nodes \ref{sec:static_node_discovery}
\begin{lstlisting}[language=bash,keywords={}]
[sms](*\#*) ochami discover static -f yaml -d @/opt/ohpc/admin/nodes/nodes.yaml
\end{lstlisting}
% ohpc_validation_newline
% end_ohpc_run



\section{Install \OHPC{} Components}\label{sec:basic_install}
\input{common/install_ohpc_components_intro}

\subsection{Enable \OHPC{} repository for local use}\label{sec:enable_repo}
% -*- mode: latex; fill-column: 120; -*-
We need to enable the OpenHPC repository on the local list of available repositories.
We then install the OPHC built applications such as the base and slurm server on the master node.
% begin_ohpc_run
% ohpc_comment_header Setup Host SLURM node
\begin{lstlisting}[language=bash,keywords={}]
[sms](*\#*) (*\install*) http://repos.openhpc.community/OpenHPC/4/EL_10/x86_64/ohpc-release-4-1.el10.x86_64.rpm
[sms](*\#*) (*\install*) dnf-plugins-core
[sms](*\#*) dnf config-manager --set-enabled crb
\end{lstlisting}
% ohpc_validation_newline
% end_ohpc_run


% begin_ohpc_run
% ohpc_validation_newline
% ohpc_validation_comment Verify OpenHPC repository has been enabled before proceeding
% ohpc_validation_newline
% ohpc_command dnf repolist | grep -q OpenHPC
% ohpc_command if [ $? -ne 0 ];then
% ohpc_command    echo "Error: OpenHPC repository must be enabled locally"
% ohpc_command    exit 1
% ohpc_command fi
% end_ohpc_run

\input{common/rocky_repos}


% begin_ohpc_run
\begin{lstlisting}[language=bash,keywords={},basicstyle=\fontencoding{T1}\fontsize{8.0}{10}\ttfamily,literate={ARCH}{\arch{}}1 {-}{-}1]
[sms](*\#*) (*\install*) epel-release

# Enable crb on sms
[sms](*\#*) /usr/bin/crb enable
\end{lstlisting}
% end_ohpc_run

Now \OHPC{} packages can be installed. To add the base package on the SMS
issue the following
% begin_ohpc_run
\begin{lstlisting}[language=bash,keywords={},basicstyle=\fontencoding{T1}\fontsize{8.0}{10}\ttfamily,literate={ARCH}{\arch{}}1 {-}{-}1]
[sms](*\#*)  (*\install*) ohpc-base
\end{lstlisting}
% end_ohpc_run

\input{common/automation}

\subsection{Setup time synchronization service on {\em master} node}\label{sec:add_ntp}
\input{common/time}


\subsection{Add resource management services on {\em master} node}\label{sec:add_rm}
\input{common/install_slurm}

\section{OpenCHAMI Image Building}\label{sec:image_building}
OpenCHAMI uses an Infrastructure as Code tool called image-builder to translate
YAML files to system images in layers to build out reusable file system layers
like container images.  This enables the rapid rebuild of images in smaller
more usable and readable layers.
OpenCHAMI used SquashFS images over S3 to serve to nodes, and Container images
through OCI.
\vspace*{0.9cm}

\subsection{Working Directory}\label{sec:image_working_directory}
% begin_ohpc_run
% ohpc_validation_newline
% ohpc_comment_header Boot Image Creation\ref{sec:image_working_directory}
% ohpc_validation_comment Creating images workspace
% ohpc_validation_newline
\begin{lstlisting}[language=bash,keywords={}]
[sms](*\#*) mkdir -p /opt/ohpc/admin/images/data
[sms](*\#*) cd /opt/ohpc/admin/images
# Register Slurm server with computes (using "configless" option)
[sms](*\#*) echo SLURMD_OPTIONS="--conf-server ${sms_ip}" > /opt/ohpc/admin/images/data/slurmd
# Handle SSH to compute nodes
[sms](*\#*) echo -e "Host ${compute_prefix}*\n\tStrictHostKeyChecking=no" >> /etc/ssh/ssh_config.d/40-ohpc.conf
\end{lstlisting}
% ohpc_validation_newline
% end_ohpc_run

\subsection{s3cmd Setup}\label{sec:s3cmd_setup}
In this recipe \OpenCHAMI{} uses {\em S3} to store and serve images and we need
to add configuration files that allow us to access and manage the {\em S3}
storage layer for {\em EFI} and boot image storage.


First lets install a CLI tool we can use to interact with the Minio service
% begin_ohpc_run
% ohpc_validation_comment Installing s3cmd CLI tool
\begin{lstlisting}[language=bash,keywords={}]
[sms](*\#*) (*\install*) s3cmd
\end{lstlisting}
% ohpc_validation_newline
% end_ohpc_run

Now we'll create a config file (maybe use a better passwd)
% begin_ohpc_run
% ohpc_validation_comment Setting S3cmd Config file
\begin{lstlisting}[language=bash,keywords={}]
[sms](*\#*) cat > ~/.s3cfg <<EOF
[sms](*\#*) # Setup endpoint
[sms](*\#*) host_base = ${sms_name}:9000
[sms](*\#*) host_bucket = ${sms_name}:9000
[sms](*\#*) bucket_location = us-east-1
[sms](*\#*) use_https = False
[sms](*\#*)
[sms](*\#*) # Setup access keys
[sms](*\#*) access_key = admin
[sms](*\#*) secret_key = password1234
[sms](*\#*)
[sms](*\#*) # Enable S3 v4 signature APIs
[sms](*\#*) signature_v2 = False
[sms](*\#*) EOF
\end{lstlisting}
% ohpc_validation_newline
% end_ohpc_run

Now we can use the tool to create the buckets we'll need and set ACLs to public
% begin_ohpc_run
% ohpc_validation_comment Creating S3 buckets
\begin{lstlisting}[language=bash,keywords={}]
[sms](*\#*) s3cmd mb s3://efi
[sms](*\#*) s3cmd setacl s3://efi --acl-public
[sms](*\#*) s3cmd mb s3://boot-images
[sms](*\#*) s3cmd setacl s3://boot-images --acl-public
\end{lstlisting}
% ohpc_validation_newline
% end_ohpc_run

The nodes need to be able to read from the S3 bucket without
authentication so we need to create some read policies
% begin_ohpc_run
% ohpc_validation_comment Setting S3 Bucket policies
\begin{lstlisting}[language=bash,keywords={}]
[sms](*\#*) cat > /opt/ohpc/admin/images/public-read-boot.json <<EOF
[sms](*\#*) {
[sms](*\#*)   "Version":"2012-10-17",
[sms](*\#*)   "Statement":[
[sms](*\#*)     {
[sms](*\#*)       "Effect":"Allow",
[sms](*\#*)      "Principal":"*",
[sms](*\#*)       "Action":["s3:GetObject"],
[sms](*\#*)       "Resource":["arn:aws:s3:::boot-images/*"]
[sms](*\#*)     }
[sms](*\#*)   ]
[sms](*\#*) }
[sms](*\#*) EOF
[sms](*\#*) cat > /opt/ohpc/admin/images/public-read-efi.json <<EOF
[sms](*\#*) {
[sms](*\#*)   "Version":"2012-10-17",
[sms](*\#*)   "Statement":[
[sms](*\#*)     {
[sms](*\#*)       "Effect":"Allow",
[sms](*\#*)       "Principal":"*",
[sms](*\#*)       "Action":["s3:GetObject"],
[sms](*\#*)       "Resource":["arn:aws:s3:::efi/*"]
[sms](*\#*)     }
[sms](*\#*)   ]
[sms](*\#*) }
[sms](*\#*) EOF
\end{lstlisting}
% ohpc_validation_newline
% end_ohpc_run

Lastly we can now apply these policies. Objects stored in the buckets should be "publicy" available
% begin_ohpc_run
% ohpc_validation_comment Applying S3 policies
\begin{lstlisting}[language=bash,keywords={}]
[sms](*\#*) s3cmd setpolicy /opt/ohpc/admin/images/public-read-boot.json s3://boot-images --host=${sms_ip}:9000 --host-bucket=${sms_ip}:9000
[sms](*\#*) s3cmd setpolicy /opt/ohpc/admin/images/public-read-efi.json s3://efi --host=${sms_ip}:9000 --host-bucket=${sms_ip}:9000
\end{lstlisting}
% ohpc_validation_newline
% end_ohpc_run


\subsection{Building the Base OS Image}\label{sec:build_base_os}
% -*- mode: latex; fill-column: 120; -*-
The \href{https://github.com/OpenCHAMI/image-builder}{Image Builder} is
designed to use layered building of system images, as shown below to enable the
use of different layers which can be modified and added to to build out
different capabilities into image layers so a change of one layer does not
require a change and rebuild of all other layers.

\begin{center}
\begin{tcolorbox}[]
\small The image-build is an optional piece of the \OpenCHAMI{} software stack.
If you have an existing image build process that creates the necessary boot
artifacts you can use that instead.
\end{tcolorbox}
\end{center}

First let's define the lowest layer of the image which disable the installation
of documentation files to reduce the size of the image:

% begin_ohpc_run
% ohpc_validation_comment Define nodocs image layer
\begin{lstlisting}[language=bash,keywords={}]
[sms](*\#*) echo -e "%_excludedocs 1" >> /opt/ohpc/admin/images/data/rpmmacros
[sms](*\#*) cat > /opt/ohpc/admin/images/nodocs.yaml <<EOF
[sms](*\#*) options:
[sms](*\#*)   layer_type: 'base'
[sms](*\#*)   name: 'rocky-nodocs'
[sms](*\#*)   publish_tags: '10'
[sms](*\#*)   pkg_manager: 'dnf'
[sms](*\#*)   parent: 'scratch'
[sms](*\#*)   publish_registry: 'ohpc-lenovo-sms:5000/c'
[sms](*\#*)   registry_opts_push:
[sms](*\#*)     - '--tls-verify=false'
[sms](*\#*) copyfiles:
[sms](*\#*)   - src: '/data/rpmmacros'
[sms](*\#*)     dest: '/root/.rpmmacros'
[sms](*\#*) EOF
\end{lstlisting}
% ohpc_validation_newline
% end_ohpc_run

Build this first layer:

% begin_ohpc_run
% ohpc_comment_header Build base layer
\begin{lstlisting}[language=bash,literate={-}{-}1,keywords={},upquote=true,keepspaces]
[sms](*\#*) podman run \
		--rm \
		--device /dev/fuse \
		--network host \
		-v /opt/ohpc/admin/images/data:/data \
		-v /opt/ohpc/admin/images/nodocs.yaml:/home/builder/config.yaml \
		ghcr.io/openchami/image-build:latest image-build \
			--config config.yaml \
			--log-level INFO
\end{lstlisting}
% ohpc_validation_newline
% end_ohpc_run

Now let's define a base layer we can use to build on. This base layer will
install common packages and also create out initramfs. In this recipe it's
created to build an initramfs that supports liveOS so the image will boot into
memory.

% begin_ohpc_run
% ohpc_validation_comment Define base image layer
\begin{lstlisting}[language=bash,keywords={}]
[sms](*\#*) cat > /opt/ohpc/admin/images/base.yaml <<EOF
[sms](*\#*) options:
[sms](*\#*)   layer_type: 'base'
[sms](*\#*)   name: 'rocky-base'
[sms](*\#*)   publish_tags: '10'
[sms](*\#*)   pkg_manager: 'dnf'
[sms](*\#*)   parent: '${sms_name}:5000/${compute_prefix}/rocky-nodocs:10'
[sms](*\#*)   publish_registry: '${sms_name}:5000/${compute_prefix}'
[sms](*\#*)   registry_opts_pull:
[sms](*\#*)     - '--tls-verify=false'
[sms](*\#*)   registry_opts_push:
[sms](*\#*)     - '--tls-verify=false'
[sms](*\#*) repos:
[sms](*\#*)   - alias: 'Rocky_10_BaseOS'
[sms](*\#*)     url: 'https://dl.rockylinux.org/pub/rocky/10/BaseOS/x86_64/os/'
[sms](*\#*)     gpg: 'https://dl.rockylinux.org/pub/rocky/RPM-GPG-KEY-Rocky-10'
[sms](*\#*)   - alias: 'Rocky_10_AppStream'
[sms](*\#*)     url: 'https://dl.rockylinux.org/pub/rocky/10/AppStream/x86_64/os/'
[sms](*\#*)     gpg: 'https://dl.rockylinux.org/pub/rocky/RPM-GPG-KEY-Rocky-10'
[sms](*\#*) package_groups:
[sms](*\#*)   - 'Minimal Install'
[sms](*\#*)   - 'Development Tools'
[sms](*\#*) packages:
[sms](*\#*)   - kernel
[sms](*\#*)   - wget
[sms](*\#*)   - dracut-live
[sms](*\#*)   - cloud-init
[sms](*\#*)   - chrony
[sms](*\#*)   - rsyslog
[sms](*\#*)   - sudo
[sms](*\#*) cmds:
[sms](*\#*)   - cmd: >
[sms](*\#*)       DRACUT_NO_XATTR=1 dracut --add "dmsquash-live livenet network-manager"
[sms](*\#*)       --kver "\$(basename /lib/modules/*)"
[sms](*\#*)       -N -f --logfile /tmp/dracut.log 2>/dev/null
[sms](*\#*)     loglevel: INFO
[sms](*\#*)   - cmd: 'echo DRACUT LOG:; cat /tmp/dracut.log'
[sms](*\#*)     loglevel: INFO
[sms](*\#*) EOF
\end{lstlisting}
% ohpc_validation_newline
% end_ohpc_run

Now lets build this base layer. This will take about 5 minutes but we should
not have to rebuild it very often.

% begin_ohpc_run
% ohpc_comment_header Build base layer
\begin{lstlisting}[language=bash,literate={-}{-}1,keywords={},upquote=true,keepspaces]
[sms](*\#*) podman run \
		--rm \
		--device /dev/fuse \
		--network host \
		-v /opt/ohpc/admin/images/base.yaml:/home/builder/config.yaml \
		ghcr.io/openchami/image-build:latest image-build \
			--config config.yaml \
			--log-level DEBUG
\end{lstlisting}
% ohpc_validation_newline
% end_ohpc_run

Now we can build on top of the base layer by setting it to be the parent of the compute-base layer

% begin_ohpc_run
% ohpc_comment_header Define the next layer
\begin{lstlisting}[language=bash,keywords={}]
[sms](*\#*) cat > /opt/ohpc/admin/images/compute-base.yaml << EOF
[sms](*\#*) options:
[sms](*\#*)   layer_type: 'base'
[sms](*\#*)   name: 'compute-base'
[sms](*\#*)   publish_tags:
[sms](*\#*)     - '10'
[sms](*\#*)   pkg_manager: 'dnf'
[sms](*\#*)   parent: '${sms_name}:5000/${compute_prefix}/rocky-base:10'
[sms](*\#*)   publish_registry: '${sms_name}:5000/${compute_prefix}'
[sms](*\#*)   registry_opts_pull:
[sms](*\#*)     - '--tls-verify=false'
[sms](*\#*)   registry_opts_push:
[sms](*\#*)     - '--tls-verify=false'
[sms](*\#*) repos:
[sms](*\#*)   - alias: 'Epel10'
[sms](*\#*)     url: 'https://dl.fedoraproject.org/pub/epel/10z/Everything/x86_64/'
[sms](*\#*)     gpg: 'https://dl.fedoraproject.org/pub/epel/RPM-GPG-KEY-EPEL-10'
[sms](*\#*) package_groups:
[sms](*\#*)   - 'InfiniBand Support'
[sms](*\#*) packages:
[sms](*\#*)   - vim
[sms](*\#*)   - nfs-utils
[sms](*\#*)   - tcpdump
[sms](*\#*)   - traceroute
[sms](*\#*)   - git
[sms](*\#*)   - http://repos.openhpc.community/OpenHPC/4/EL_10/x86_64/ohpc-release-4-1.el10.x86_64.rpm
[sms](*\#*) cmds:
[sms](*\#*)   - cmd: dnf config-manager --set-enabled crb
[sms](*\#*)     loglevel: INFO
[sms](*\#*) EOF
\end{lstlisting}
% ohpc_validation_newline
% end_ohpc_run

Then we can build it. Any modifications to the above configuration will only
require a rebuild of this layer

% begin_ohpc_run
% ohpc_comment_header Build the next layer
\begin{lstlisting}[language=bash,literate={-}{-}1,keywords={},upquote=true,keepspaces]
[sms](*\#*) podman run \
                --rm \
                --device /dev/fuse \
                --network host \
                -v /opt/ohpc/admin/images/compute-base.yaml:/home/builder/config.yaml \
                ghcr.io/openchami/image-build:latest image-build \
                        --config config.yaml \
                        --log-level DEBUG
\end{lstlisting}
% ohpc_validation_newline
% end_ohpc_run

You can keep adding layers as you see fit, and branch off if so desired. We
will use the next layer to add additional packages later on so no need to build
it just yet.

% begin_ohpc_run
% ohpc_validation_comment Define a production compute layer
\begin{lstlisting}[language=bash,keywords={}]
[sms](*\#*) cp /etc/hosts /opt/ohpc/admin/images/data
[sms](*\#*) cat > /opt/ohpc/admin/images/compute-prod.yaml << EOF
[sms](*\#*) options:
[sms](*\#*)   layer_type: 'base'
[sms](*\#*)   name: 'compute-prod'
[sms](*\#*)   publish_tags:
[sms](*\#*)     - '10'
[sms](*\#*)   pkg_manager: 'dnf'
[sms](*\#*)   parent: '${sms_name}:5000/${compute_prefix}/compute-base:10'
[sms](*\#*)   publish_registry: '${sms_name}:5000/${compute_prefix}'
[sms](*\#*)   registry_opts_pull:
[sms](*\#*)     - '--tls-verify=false'
[sms](*\#*)   registry_opts_push:
[sms](*\#*)     - '--tls-verify=false'
[sms](*\#*)   # Publish SquashFS image to local S3
[sms](*\#*)   publish_s3: 'http://${sms_name}:9000'
[sms](*\#*)   s3_prefix: 'compute/base/'
[sms](*\#*)   s3_bucket: 'boot-images'
[sms](*\#*)
[sms](*\#*)   # Publish OCI image to container registry
[sms](*\#*)   #
[sms](*\#*)   # This is the only way to be able to reuse this image as
[sms](*\#*)   # a parent for another image layer.
[sms](*\#*)   publish_registry: '${sms_name}:5000/${compute_prefix}'
[sms](*\#*)   registry_opts_push:
[sms](*\#*)     - '--tls-verify=false'
[sms](*\#*) copyfiles:
[sms](*\#*)   - src: '/data/hosts'
[sms](*\#*)     dest: '/etc/hosts'
[sms](*\#*)   - src: '/data/slurmd'
[sms](*\#*)     dest: '/etc/sysconfig/slurmd'
[sms](*\#*) packages:
[sms](*\#*)   - ohpc-base-compute
[sms](*\#*)   - ohpc-slurm-client
[sms](*\#*)   - pdsh-ohpc
[sms](*\#*)   - lmod-ohpc
[sms](*\#*) cmds:
[sms](*\#*)   - cmd: systemctl disable firewalld
[sms](*\#*)     loglevel: INFO
[sms](*\#*)   - cmd: echo '${sms_ip}:/home /home nfs nfsvers=3,nodev,nosuid 0 0' >> /etc/fstab
[sms](*\#*)     loglevel: INFO
[sms](*\#*)   - cmd: echo '${sms_ip}:/opt/ohpc/pub /opt/ohpc/pub nfs nfsvers=3,nodev,nosuid 0 0' >> /etc/fstab
[sms](*\#*)     loglevel: INFO
[sms](*\#*)   - cmd: mkdir -p /opt/ohpc/pub
[sms](*\#*)     loglevel: INFO
[sms](*\#*)   - cmd: echo "account required pam_slurm.so" >> /etc/pam.d/sshd
[sms](*\#*)     loglevel: INFO
[sms](*\#*)   - cmd: systemctl enable munge slurmd
[sms](*\#*)     loglevel: INFO
[sms](*\#*)   - cmd: ln -sf ../usr/share/zoneinfo/UTC /etc/localtime
[sms](*\#*)     loglevel: INFO
[sms](*\#*) EOF
\end{lstlisting}
% ohpc_validation_newline
% end_ohpc_run


\vspace*{-0.25cm}
\subsubsection{Customize system configuration} \label{sec:master_customization}
Here we set up \NFS{}  mounting of a
\$HOME file system and the public \OHPC{} install path (\texttt{/opt/ohpc/pub})
that will be hosted by the {\em SMS} host in this  example configuration.

\vspace*{0.15cm}
% begin_ohpc_run
% ohpc_comment_header Customize system configuration \ref{sec:master_customization}
\begin{lstlisting}[language=bash,literate={-}{-}1,keywords={},upquote=true]
# Export /home and OpenHPC public packages from master server
[sms](*\#*) echo "/home *(rw,no_subtree_check,fsid=10,no_root_squash)" >> /etc/exports
[sms](*\#*) echo "/opt/ohpc/pub *(ro,no_subtree_check,fsid=11)" >> /etc/exports
[sms](*\#*) exportfs -a
[sms](*\#*) systemctl restart nfs-server
[sms](*\#*) systemctl enable nfs-server
\end{lstlisting}
% end_ohpc_run



% Additional commands when additional computes are requested

% begin_ohpc_run
% ohpc_validation_newline
% ohpc_validation_comment Update basic slurm configuration if additional computes defined
% ohpc_validation_comment This is performed on the SMS, nodes will pick it up config file is copied there later
% ohpc_command if [ ${num_computes} -gt 4 ];then
% ohpc_command    perl -pi -e "s/^NodeName=(\S+)/NodeName=${compute_prefix}[1-${num_computes}]/" /etc/slurm/slurm.conf
% ohpc_command    perl -pi -e "s/^PartitionName=normal Nodes=(\S+)/PartitionName=normal Nodes=${compute_prefix}[1-${num_computes}]/" /etc/slurm/slurm.conf
% ohpc_command fi
% end_ohpc_run

%\clearpage

\paragraph{Increase locked memory limits}
In order to utilize \InfiniBand{} or Omni-Path as the underlying high speed
interconnect, it is generally necessary to increase the locked memory settings
for system users. This can be accomplished by adding the
\texttt{/etc/security/limits.d/40-ohpc-limits.conf} file and this should be
performed on all job submission hosts. In this recipe, jobs are submitted from
the {\em SMS} host, and the following commands can be used to update the
maximum locked memory settings on both the {\em SMS} host and compute nodes:

% begin_ohpc_run ohpc_validation_newline ohpc_validation_comment Update memlock
% settings
\begin{lstlisting}[language=bash,literate={-}{-}1,keywords={},upquote=true,keepspaces]
# Update memlock settings on SMS
[sms](*\#*) echo '* soft memlock unlimited' >> /etc/security/limits.d/40-ohpc-limits.conf
[sms](*\#*) echo '* hard memlock unlimited' >> /etc/security/limits.d/40-ohpc-limits.conf
# Update memlock settings on compute
[sms](*\#*) yq -i \
	'.cmds += [{"cmd":"echo \"* soft memlock unlimited\" >> /etc/security/limits.d/40-ohpc-limits.conf"}]' \
	/opt/ohpc/admin/images/compute-prod.yaml
[sms](*\#*) yq -i \
	'.cmds += [{"cmd":"echo \"* hard memlock unlimited\" >> /etc/security/limits.d/40-ohpc-limits.conf"}]' \
	/opt/ohpc/admin/images/compute-prod.yaml
\end{lstlisting}
% end_ohpc_run


\paragraph{Add \Lustre{} client}\label{sec:lustre_client}
\input{common/lustre-client}
% begin_ohpc_run
% ohpc_validation_newline
% ohpc_validation_comment Enable Optional packages
% ohpc_validation_newline
% ohpc_command if [[ ${enable_lustre_client} -eq 1 ]];then
% ohpc_indent 5

% ohpc_validation_comment Install Lustre client on master
\begin{lstlisting}[language=bash,keywords={},upquote=true]
# Add Lustre client software to master host
[sms](*\#*) (*\install*) lustre-client-ohpc
\end{lstlisting}
% end_ohpc_run

% begin_ohpc_run
% ohpc_validation_comment Enable lustre in WW compute image
\begin{lstlisting}[language=bash,keywords={},upquote=true]
# Include Lustre client software in compute image
[sms](*\#*) yq -i '.packages += ["lustre-client-ohpc"]' /opt/ohpc/admin/images/compute-prod.yaml

\end{lstlisting}
% end_ohpc_run

\input{common/lustre-client-post-openchami}

\vspace*{0.4cm}

\paragraph{Add GPU driver}
% begin_ohpc_run
\noindent If planning to install the NVIDIA toolkit, register the
following additional path (\texttt{/opt/nvidia}) to share with the compute nodes:
% ohpc_command if [[ ${enable_nvidia_gpu_driver} -eq 1 ]];then
% ohpc_indent 5
\begin{lstlisting}[language=bash,keywords={},upquote=true,keepspaces]
# Add NVIDIA GPU driver repository to the SMS
[sms](*\#*) (*\install*) cuda-repo-ohpc
# Add NVIDIA GPU driver repository to the compute nodes
[sms](*\#*) yq -i '.repos += [{"alias":"cuda-rhel9-x86_64", "url":"https://developer.download.nvidia.com/compute/cuda/repos/rhel9/x86_64", "gpg":"https://developer.download.nvidia.com/compute/cuda/repos/rhel9/x86_64/D42D0685.pub" }]' /opt/ohpc/admin/images/compute-prod.yaml
# Install the GPU driver on the compute nodes
[sms](*\#*) yq -i '.modules += {"install": ["nvidia-driver:latest-dkms"]}' /opt/ohpc/admin/images/compute-prod.yaml
# Enable DKMS service to automatically rebuild driver
[sms](*\#*) yq -i '.cmds += [{"cmd":"systemctl enable dkms"}]' /opt/ohpc/admin/images/compute-prod.yaml
# Install the toolkit on the SMS
[sms](*\#*) (*\install*) cuda-devel-ohpc nvidia-driver-cuda
# (Optional) Setup NFS mount for /opt/nvidia
[sms](*\#*) echo "/opt/nvidia *(ro,no_subtree_check)" >> /etc/exports
[sms](*\#*) exportfs -a

# Create mount point and file system mount
[sms](*\#*) nodeshell compute mkdir /opt/nvidia
[sms](*\#*) nodeshell compute echo \
        "\""${sms_ip}:/opt/nvidia /opt/nvidia nfs nfsvers=3,nodev,nosuid 0 0"\"" \>\> /etc/fstab
# Mount NFS shares
[sms](*\#*) nodeshell compute mount /opt/nvidia

\end{lstlisting}
% ohpc_indent 0
% ohpc_command fi


\paragraph{Add \clustershell{}}
\input{common/clustershell}

\paragraph{Add \genders{}}
\input{common/genders}

\paragraph{Add Magpie}
\input{common/magpie}


\paragraph{Add \conman{}} \label{sec:add_conman}
\input{common/conman}

\paragraph{Add \nhc{}} \label{sec:add_nhc}
\input{common/nhc_openchami}
\input{common/nhc_slurm}


\paragraph{Build Production Compute Image}
% -*- mode: latex; fill-column: 120; -*-
Rebuild Final Compute node image with all optional packages

% begin_ohpc_run
% ohpc_comment_header Build Compute Image
\begin{lstlisting}[language=bash,literate={-}{-}1,keywords={},upquote=true,keepspaces]
[sms](*\#*) podman run \
		--rm \
		--device /dev/fuse \
		--network host \
		-e S3_ACCESS=admin \
		-e S3_SECRET=password1234 \
		-v /opt/ohpc/admin/images/data:/data \
		-v /opt/ohpc/admin/images/compute-prod.yaml:/home/builder/config.yaml \
		ghcr.io/openchami/image-build:latest \
			image-build \
			--config config.yaml
\end{lstlisting}
% ohpc_validation_newline
% end_ohpc_run


\subsection{Setting Boot Parameters}\label{sec:boot_params_set}
% -*- mode: latex; fill-column: 120; -*-
Boot Script Service (BSS) is what provides the path to vmlinuz and the initrd, and provides the kernel parameters to the node during iPXE.

Get some data from our images to create our BSS payload.
% begin_ohpc_run
% ohpc_comment_header Get info from image to use for BSS \ref{sec:boot_params_set}
\begin{lstlisting}[language=bash,keywords={}]
[sms](*\#*) export s3Output=$(s3cmd ls -Hr s3://boot-images/)
[sms](*\#*) export os=$(echo "$s3Output" | awk '{print $4}' | grep 'boot-images/compute')
[sms](*\#*) export ramfs=$(echo "$s3Output" | awk '{print $4}' | grep 'initram')
[sms](*\#*) export hdrs=$(echo "$s3Output" | awk '{print $4}' | grep 'vmlinuz')
[sms](*\#*) export params="nomodeset ro root=live:http://${sms_ip}:9000/${os:5} \
                           ip=dhcp overlayroot=tmpfs overlayroot_cfgdisk=disabled \
			   apparmor=0 selinux=0 console=ttyS0,115200 ip6=off \
			   cloud-init=enabled ds=nocloud-net;s=http://${sms_ip}:8081/cloud-init"
\end{lstlisting}
% ohpc_validation_newline
% end_ohpc_run

Create a yaml payload for BSS
% begin_ohpc_run
% ohpc_comment_header Define the BSS payload \ref{sec:boot_params_set}
\begin{lstlisting}[language=bash,keywords={}]
[sms](*\#*) cat >> /opt/ohpc/admin/nodes/compute-boot.yaml <<EOF
[sms](*\#*) kernel: 'http://${sms_ip}:9000/${hdrs:5}'
[sms](*\#*) initrd: 'http://${sms_ip}:9000/${ramfs:5}'
[sms](*\#*) params: '${params}'
[sms](*\#*) macs:
[sms](*\#*) EOF
[sms](*\#*) for((i=0; i < $num_computes; i++)); do
                 echo "  - ${c_mac[$i]}" >> /opt/ohpc/admin/nodes/compute-boot.yaml
      done
\end{lstlisting}
% ohpc_validation_newline
% end_ohpc_run

And then finally upload our payload
% begin_ohpc_run
% ohpc_comment_header Upload the BSS payload \ref{sec:boot_params_set}
\begin{lstlisting}[language=bash,keywords={}]
[sms](*\#*) ochami bss boot params set -f yaml -d @/opt/ohpc/admin/nodes/compute-boot.yaml
\end{lstlisting}
% ohpc_validation_newline
% end_ohpc_run


\subsection{Building CloudInit Image}\label{sec:cloudinit}
% -*- mode: latex; fill-column: 120; -*-
\OpenCHAMI{} uses cloudinit for post boot node configuration. We will establish
some basic cloudinit configurations here. But later we will build a cloudinit
template that injects the munge key into the system at boot.


Let's create a directory to hold our default cloud-init configuration and create an
SSH key pair that we can use for SSH access later

% begin_ohpc_run
% ohpc_comment_header Create SSH keys \ref{sec:cloudinit}
\begin{lstlisting}[language=bash,keywords={}]
[sms](*\#*) mkdir -p /opt/ohpc/admin/cloud-init
[sms](*\#*) cd /opt/ohpc/admin/cloud-init
[sms](*\#*) ssh-keygen -t ed25519 -q -f "$HOME/.ssh/id_openchami" -N ""
[sms](*\#*) export master_key=$(cat ~/.ssh/id_openchami.pub)
\end{lstlisting}
% ohpc_validation_newline
% end_ohpc_run

Now we can define the payload to update the defaults
% begin_ohpc_run
% ohpc_comment_header Create cloud-init payload \ref{sec:cloudinit}
\begin{lstlisting}[language=bash,keywords={},upquote=true,keepspaces]
[sms](*\#*) cat > /opt/ohpc/admin/cloud-init/defaults.yaml <<EOF
[sms](*\#*) ---
[sms](*\#*) base-url: "http://${sms_ip}:8081/cloud-init"
[sms](*\#*) cluster-name: "${compute_prefix^^}"
[sms](*\#*) nid-length: 1
[sms](*\#*) public-keys:
[sms](*\#*) - "${master_key}"
[sms](*\#*) short-name: "${compute_prefix}"
[sms](*\#*) EOF
[sms](*\#*) for i in /root/.ssh/id*pub; do \
		file=${i} yq '.public-keys +=loadstr(strenv(file))' -i /opt/ohpc/admin/cloud-init/defaults.yaml; \
	done
\end{lstlisting}
% ohpc_validation_newline
% end_ohpc_run

And then finally upload our payload
% begin_ohpc_run
% ohpc_comment_header Upload the cloud-init defaults \ref{sec:boot_params_set}
\begin{lstlisting}[language=bash,keywords={}]
[sms](*\#*) ochami cloud-init defaults set -f yaml -d @/opt/ohpc/admin/cloud-init/defaults.yaml
\end{lstlisting}
% ohpc_validation_newline
% end_ohpc_run



%\subsubsection{Identify files for synchronization} \label{sec:file_import}
%\input{common/import_confluent_files}
%\input{common/import_confluent_files_slurm}

%%%\subsubsection{Optional kernel arguments} \label{sec:optional_kargs}
%%%\input{common/conman_post}

\section{Munge Key Injection} \label{sec:cloudinit-munge}
\OpenCHAMI{} provides a cloud-init server that can be used, in conjunction with
the cloud-init client, to provide post-boot configurations. The configuration
below is a minimal configuration that should at least let the root user login
using the public key and distribute the munge key file to all compute nodes.

% begin_ohpc_run ohpc_comment_header Cloudinit \ref{sec:cloudinit}
\begin{lstlisting}[language=bash,keywords={}]
[sms](*\#*) export munge_key=$(cat /etc/munge/munge.key | base64 -w 0)
[sms](*\#*) cat > /opt/ohpc/admin/cloud-init/compute.yaml <<EOF
[sms](*\#*) - name: compute
[sms](*\#*)   description: "compute template"
[sms](*\#*)   file:
[sms](*\#*)     encoding: plain
[sms](*\#*)     content: |
[sms](*\#*)       ## template: jinja
[sms](*\#*)       #cloud-config
[sms](*\#*)       merge_how:
[sms](*\#*)       - name: list
[sms](*\#*)         settings: [append]
[sms](*\#*)       - name: dict
[sms](*\#*)         settings: [no_replace, recurse_list]
[sms](*\#*)       users:
[sms](*\#*)         - name: root
[sms](*\#*)           ssh_authorized_keys: {{ ds.meta_data.instance_data.v1.public_keys }}
[sms](*\#*)       write_files:
[sms](*\#*)         - content: ${munge_key}
[sms](*\#*)           path: /etc/munge/munge.key
[sms](*\#*)           permissions: '0400'
[sms](*\#*)           owner: munge:munge
[sms](*\#*)           encoding: base64
[sms](*\#*) EOF
\end{lstlisting}
% ohpc_validation_newline
% end_ohpc_run

Once you are done updating cloud-init you can post the file to the cloud-init server
% begin_ohpc_run
\begin{lstlisting}[language=bash,keywords={}]
[sms](*\#*) ochami cloud-init group set -f yaml -d @/opt/ohpc/admin/cloud-init/compute.yaml
\end{lstlisting}
% ohpc_validation_newline
% end_ohpc_run


%\vspace*{-0.25cm}
\subsection{Boot compute nodes} \label{sec:boot_computes}
\input{common/reset_computes}

\section{Install \OHPC{} Development Components}
\input{common/dev_intro}

%\vspace*{-0.15cm}
%\clearpage
\subsection{Development Tools} \label{sec:install_dev_tools}
\input{common/dev_tools}

\vspace*{-0.15cm}
\subsection{Compilers} \label{sec:install_compilers}
\input{common/compilers}

%\clearpage
\subsection{MPI Stacks} \label{sec:mpi}
\input{common/mpi_slurm}

\subsection{Performance Tools} \label{sec:install_perf_tools}
\input{common/perf_tools}

\subsection{Setup default development environment}
\input{common/default_dev}

%\vspace*{0.2cm}
\subsection{3rd Party Libraries and Tools} \label{sec:3rdparty}
\input{common/third_party_libs_intro}

\input{common/third_party_libs}
\input{common/third_party_mpi_libs_x86}

\vspace*{.6cm}
\subsection{Optional Development Tool Builds} \label{sec:3rdparty_intel}
\input{common/oneapi_enabled_builds_slurm}

\section{Resource Manager Startup} \label{sec:rms_startup}
\input{common/slurm_startup}

\clearpage
\appendix
{\bf \LARGE \centerline{Appendices}} \vspace*{0.2cm}

\addcontentsline{toc}{section}{Appendices}
\renewcommand{\thesubsection}{\Alph{subsection}}

\input{common/automation_appendix}
\input{common/test_suite}
\input{common/customization_appendix_centos}
\input{manifest}
\clearpage
\subsection{Package Signatures}
%\addcontentsline{toc}{section}{Appendix B - Package Signatures}

All of the RPMs provided via the \OHPC{} repository are signed with a GPG
signature. By default, the underlying package managers will verify these signatures during
installation to ensure that packages have not been altered. The RPMs can also
be manually verified and the public signing key fingerprint for the latest
repository is shown below: \\

\texttt{Fingerprint: 5E33 3CA3 A1BD BBC9 DF14 9D74 09AD FAE4 {\bf D722A692}} \\

\noindent The following command can be used to verify an RPM once it
has been downloaded locally by confirming if the package is signed, and if so,
indicating which key was used to sign it. The example below highlights usage
for a local copy of the \texttt{docs-ohpc} package and illustrates how the {\em
key ID} matches the fingerprint shown above.

\begin{lstlisting}[language=bash,keywords={}]
[sms](*\#*) rpm --checksig -v ohpc-release-4-1.el10.x86_64.rpm
ohpc-release-4-1.el10.x86_64.rpm:
    Header V3 RSA/SHA256 Signature, key ID d722a692: OK
    Header SHA256 digest: OK
    Header SHA1 digest: OK
    Payload SHA256 digest: OK
    V3 RSA/SHA256 Signature, key ID d722a692: OK
    MD5 digest: OK

\end{lstlisting}


\end{document}
